\documentclass[aspectratio=169]{beamer}

% =====================================================================
% ART DECK 04 -- MONOCHROME. Silver light on a black void.
% A fresh, comprehensive outline of the whole project, rendered as
% black-and-white fine-art TikZ: no hue, the full tonal range
% (white -> silver -> ash -> void) and pure form carry everything.
% Brightness encodes meaning. Two pdflatex passes (overlay).
% =====================================================================

\usepackage{tikz}
\usepackage[T1]{fontenc}
\usetikzlibrary{shadings, fadings, shapes.geometric, shapes.misc,
  decorations.pathmorphing, decorations.pathreplacing, calc, backgrounds,
  positioning, arrows.meta, shapes.symbols}

\usetheme{default}
\setbeamertemplate{navigation symbols}{}
\setbeamercolor{background canvas}{bg=black}
\setbeamertemplate{footline}{}

% ---- monochrome palette (a hair cool, for silver elegance) ----------------
\definecolor{void}{HTML}{070809}
\definecolor{ink}{HTML}{0B0C0E}
\definecolor{slate}{HTML}{343B44}
\definecolor{ash}{HTML}{767F8A}
\definecolor{smoke}{HTML}{B4BCC5}
\definecolor{silver}{HTML}{DCE2E9}
\definecolor{paper}{HTML}{F4F7FA}

\tikzfading[name=rglow,  inner color=transparent!0,  outer color=transparent!100]
\tikzfading[name=fadeup, top color=transparent!0,    bottom color=transparent!100]

\pgfmathsetseed{104}

\newcommand{\art}[1]{%
  \begin{frame}[plain]
  \begin{tikzpicture}[remember picture, overlay,
      shift={(current page.south west)}, x=1cm, y=1cm]
    #1
  \end{tikzpicture}
  \end{frame}}

\newcommand{\halo}[4]{\fill[#3, opacity=#4, path fading=rglow] #1 circle (#2);}

% luminous star with lens-flare cross (core always white)
\newcommand{\spark}[3]{%
  \fill[#3, opacity=0.50, path fading=rglow] #1 circle ({#2*7});
  \fill[#3, opacity=0.85, path fading=rglow] #1 circle ({#2*2.8});
  \begin{scope}[blend mode=screen]
    \fill[white, opacity=0.55, path fading=rglow] #1 ellipse ({#2*9} and {#2*0.45});
    \fill[white, opacity=0.55, path fading=rglow] #1 ellipse ({#2*0.45} and {#2*9});
  \end{scope}
  \fill[white] #1 circle (#2);}

% glowing silver wire (triple stroke)
\newcommand{\neon}[3]{%
  \draw[#3, line width=7pt, opacity=0.08] #1 -- #2;
  \draw[#3, line width=3pt, opacity=0.26] #1 -- #2;
  \draw[white, line width=0.7pt, opacity=0.80] #1 -- #2;}

% white star dust
\newcommand{\dust}[1]{%
  \foreach \i in {1,...,#1}{%
     \pgfmathsetmacro{\sx}{rnd*16}\pgfmathsetmacro{\sy}{rnd*9}%
     \pgfmathsetmacro{\sr}{0.005+rnd*0.020}\pgfmathsetmacro{\so}{0.08+rnd*0.50}%
     \fill[white, opacity=\so] (\sx,\sy) circle (\sr);}}

% soft silver bokeh
\newcommand{\bokehm}[1]{%
  \begin{scope}[blend mode=screen]
  \foreach \i in {1,...,#1}{%
     \pgfmathsetmacro{\bx}{rnd*16}\pgfmathsetmacro{\by}{rnd*9}%
     \pgfmathsetmacro{\br}{0.05+rnd*0.45}\pgfmathsetmacro{\bo}{0.05+rnd*0.20}%
     \pgfmathtruncatemacro{\ci}{rnd*3.999}%
     \ifcase\ci \def\bc{silver}\or\def\bc{smoke}\or\def\bc{ash}\or\def\bc{white}\fi
     \fill[\bc, opacity=\bo, path fading=rglow] (\bx,\by) circle (\br);}%
  \end{scope}}

% volumetric rays
\newcommand{\rays}[3]{%
  \begin{scope}[blend mode=screen]
  \foreach \k in {1,...,#3}{%
     \pgfmathsetmacro{\ang}{rnd*360}\pgfmathsetmacro{\len}{4.5+rnd*7}%
     \pgfmathsetmacro{\wd}{0.4+rnd*1.4}%
     \fill[#2, opacity=0.05, path fading=rglow]
        #1 -- ($#1+(\ang-\wd:\len)$) -- ($#1+(\ang+\wd:\len)$) -- cycle;}
  \end{scope}}

% glow halo behind text
\newcommand{\glowtext}[3]{%
  \begin{scope}[blend mode=screen]
     \fill[#2, opacity=0.45, path fading=rglow] #1 ellipse (5.2 and 1.5);
  \end{scope}
  \node at #1 {#3};}

% silver mist backdrop
\newcommand{\mist}{%
  \begin{scope}[blend mode=screen]
     \fill[silver, opacity=0.15, path fading=rglow] (5.5,6) circle (6.5);
     \fill[smoke,  opacity=0.12, path fading=rglow] (11,5.4) circle (5);
     \fill[ash,    opacity=0.22, path fading=rglow] (8,4) circle (8.5);
  \end{scope}}

% faint concentric engraving rings: {(c)}{count}{step}{color}
\newcommand{\rings}[4]{%
  \foreach \r in {1,...,#2}{%
     \draw[#4, opacity={0.16-0.014*\r}, line width=0.4pt] #1 circle ({\r*#3});}}

\begin{document}

% =====================================================================
% 1. TITLE
% =====================================================================
\art{
  \fill[void] (0,0) rectangle (16,9);
  \mist \dust{230}\bokehm{55}
  \rings{(8,5)}{6}{1.15}{ash}
  \coordinate (n1) at (3.6,7.0); \coordinate (n2) at (6.6,7.9);
  \coordinate (n3) at (10.2,7.1);\coordinate (n4) at (2.8,4.9);
  \coordinate (n5) at (6.2,5.3); \coordinate (n6) at (10.9,5.1);
  \coordinate (n7) at (13.0,6.3);\coordinate (n8) at (8.3,3.7);
  \coordinate (n9) at (4.9,3.2);
  \foreach \a/\b in {n1/n2,n2/n3,n1/n4,n4/n5,n5/n6,n3/n6,n2/n5,n6/n7,
     n5/n8,n8/n9,n4/n9,n9/n5}{\neon{(\a)}{(\b)}{silver}}
  \foreach \n/\s in {n1/0.075,n2/0.06,n3/0.07,n4/0.055,n5/0.085,
     n6/0.06,n7/0.05,n8/0.06,n9/0.07}{\spark{(\n)}{\s}{silver}}
  \begin{scope}[blend mode=screen]
     \fill[silver, opacity=0.35, path fading=rglow] (5.7,1.85) ellipse (6 and 1.5);
  \end{scope}
  \node[anchor=west, text=white] at (1.3,2.15)
    {\fontsize{48}{52}\selectfont\textbf{Rotor} {\color{ash}\textit{in}} \textbf{Rust}};
  \node[anchor=west, text=silver, opacity=0.92] at (1.4,1.2)
    {\large A machine that asks one question of every input.};
  \node[anchor=west, text=ash] at (1.4,0.5)
    {\footnotesize Jasmin Begi\'{c}\; \textbar\; Daniel Wassie};
}

% =====================================================================
% 2. THE PROBLEM
% =====================================================================
\art{
  \fill[void] (0,0) rectangle (16,9);
  \begin{scope}[blend mode=screen]
     \fill[ash, opacity=0.35, path fading=rglow] (8,4.5) circle (9);\end{scope}
  % field of dim inputs
  \foreach \x in {0,...,37}{\foreach \y in {0,...,16}{
     \pgfmathsetmacro{\jx}{\x*0.40+0.6+rnd*0.05}
     \pgfmathsetmacro{\jy}{\y*0.40+1.5+rnd*0.05}
     \fill[ash, opacity={0.22+rnd*0.28}] (\jx,\jy) circle (0.045);}}
  \dust{90}
  % the one that crashes
  \rays{(11.4,5.8)}{white}{24}\spark{(11.4,5.8)}{0.13}{white}
  \node[text=white] at (8,8.4) {\Large you cannot test every input --- the one that breaks \textit{hides}};
  \glowtext{(8,0.85)}{silver}{\fontsize{34}{36}\selectfont\textbf{\textcolor{white}{$2^{64}$} \normalsize\textcolor{silver}{possibilities, one of them fatal}}}
}

% =====================================================================
% 3. THE PIPELINE
% =====================================================================
\art{
  \fill[void] (0,0) rectangle (16,9);
  \mist \dust{60}
  \def\sy{4.6}
  \draw[silver, line width=9pt, opacity=0.07] (2.2,\sy) -- (13.4,\sy);
  \draw[smoke,  line width=3pt, opacity=0.22] (2.2,\sy) -- (13.4,\sy);
  \foreach \x/\lbl/\b in {2.2/binary/0, 6.0/model/1, 9.8/solver/0, 13.4/answer/1}{
     \fill[silver, opacity={0.30+0.25*\b}, path fading=rglow] (\x,\sy) circle (1.6);
     \shade[inner color=white, outer color=slate] (\x,\sy) circle (0.72);
     \draw[white, line width={0.8pt+0.8pt*\b}, opacity={0.5+0.4*\b}] (\x,\sy) circle (0.72);
     \node[text=ink, font=\footnotesize\bfseries] at (\x,\sy) {\lbl};}
  \foreach \m in {4.1,7.9,11.6}{\begin{scope}[blend mode=screen]
     \fill[white, opacity=0.8, path fading=rglow] (\m,\sy) circle (0.14);\end{scope}}
  \node[text=white] at (8,8.0) {\Large what we built sits in the \textbf{middle} of a pipeline};
  \node[text=silver, font=\footnotesize] at (6.0,2.6) {rotor};
  \node[text=silver, font=\footnotesize] at (13.4,2.6){visualizer};
  \node[text=ash, font=\footnotesize] at (8,1.5) {a compiled binary in $\rightarrow$ a readable verdict out};
}

% =====================================================================
% 4. WHAT A MODEL IS  (the 24-property aperture)
% =====================================================================
\art{
  \fill[void] (0,0) rectangle (16,9);
  \begin{scope}[blend mode=screen]
     \fill[ash, opacity=0.30, path fading=rglow] (8,4.7) circle (7);\end{scope}
  \dust{60}
  % aperture ring with 24 ticks (one highlighted)
  \draw[silver, opacity=0.65, line width=1pt] (8,4.7) circle (2.5);
  \draw[ash, opacity=0.4, line width=0.5pt] (8,4.7) circle (2.05);
  \foreach \i in {0,...,23}{
     \pgfmathsetmacro{\a}{90-\i*15}
     \ifnum\i=7
        \draw[white, line width=1.6pt] ({8+2.05*cos(\a)},{4.7+2.05*sin(\a)})
           -- ({8+2.85*cos(\a)},{4.7+2.85*sin(\a)});
        \spark{({8+2.95*cos(\a)},{4.7+2.95*sin(\a)})}{0.07}{white}
     \else
        \draw[silver, opacity=0.7, line width=1pt]
           ({8+2.5*cos(\a)},{4.7+2.5*sin(\a)}) -- ({8+2.75*cos(\a)},{4.7+2.75*sin(\a)});
     \fi}
  \node[text=white, font=\large] at (8,5.0) {a model};
  \node[text=silver, font=\footnotesize] at (8,4.4) {state $\cdot$ rules $\cdot$ 24 limits};
  \node[text=white] at (8,8.2) {\Large a board game: a position, the legal moves, the forbidden squares};
  \node[text=ash, font=\footnotesize] at (8,1.1)
     {24 things that must never happen --- one of them is \textit{division-by-zero}};
}

% =====================================================================
% 5. ROTOR IN RUST -- the rewrite
% =====================================================================
\art{
  \fill[void] (0,0) rectangle (16,9);
  \mist \dust{50}
  \node[text=white] at (8,8.2) {\Large one 14{,}000-line monolith \;$\rightarrow$\; four clean parts};
  \foreach \x/\y/\name/\sub in {%
      4.6/5.5/btor2/{the node graph},
      11.4/5.5/riscv/{decode the binary},
      4.6/2.9/machine/{state \& memory},
      11.4/2.9/model/{logic \& properties}}{
     \begin{scope}[blend mode=screen]
        \fill[silver, opacity=0.16, path fading=rglow] (\x,\y) circle (2.2);\end{scope}
     \draw[silver, opacity=0.8, line width=1pt, rounded corners=6pt]
        (\x-2.3,\y-0.85) rectangle (\x+2.3,\y+0.85);
     \node[text=white, font=\large\bfseries] at (\x,\y+0.18) {\name};
     \node[text=ash, font=\footnotesize] at (\x,\y-0.42) {\sub};}
  \draw[ash, opacity=0.5, line width=0.5pt] (6.9,5.5) -- (9.1,5.5);
  \draw[ash, opacity=0.5, line width=0.5pt] (6.9,2.9) -- (9.1,2.9);
  \draw[ash, opacity=0.5, line width=0.5pt] (4.6,4.65) -- (4.6,3.75);
  \draw[ash, opacity=0.5, line width=0.5pt] (11.4,4.65) -- (11.4,3.75);
  \node[text=silver, font=\footnotesize] at (8,0.9) {typed, modular, no globals --- a base built to \textit{extend}};
}

% =====================================================================
% 6. SPEED -- a comet
% =====================================================================
\art{
  \fill[void] (0,0) rectangle (16,9);
  \begin{scope}[blend mode=screen]
     \fill[ash, opacity=0.30, path fading=rglow] (8,4.5) circle (10);
     \fill[silver, opacity=0.20, path fading=rglow] (13.4,4.9) circle (4.5);\end{scope}
  \dust{110}
  \begin{scope}[blend mode=screen]
  \foreach \i in {1,...,40}{
     \pgfmathsetmacro{\my}{1.2+rnd*7.2}\pgfmathsetmacro{\mx}{1+rnd*9}
     \pgfmathsetmacro{\ml}{0.4+rnd*2.2}\pgfmathsetmacro{\mo}{0.05+rnd*0.16}
     \fill[white, opacity=\mo, path fading=rglow] (\mx,\my) ellipse ({\ml} and 0.02);}
  \end{scope}
  \begin{scope}[blend mode=screen]
     \fill[smoke, opacity=0.50, path fading=rglow]
        (1.5,2.7) .. controls (6,3.4) and (10,4.2) .. (13.4,4.9)
        .. controls (10,3.7) and (6,3.0) .. (1.5,2.5) -- cycle;
     \fill[white, opacity=0.45, path fading=rglow]
        (4.5,3.3) .. controls (8,3.9) and (11,4.4) .. (13.4,4.9)
        .. controls (11,4.1) and (8,3.6) .. (4.5,3.1) -- cycle;
  \end{scope}
  \rays{(13.5,4.9)}{white}{22}\spark{(13.5,4.9)}{0.17}{white}
  \node[text=white] at (5.2,7.2) {\fontsize{72}{74}\selectfont\textbf{1000$\times$}};
  \node[text=silver, opacity=0.92] at (5.4,5.7) {\large faster --- and a jump that big deserves suspicion};
  \node[text=ash] at (4.2,1.3)
    {\footnotesize 139\,s $\rightarrow$ 0.1\,s \qquad 428\,MB $\rightarrow$ 20\,MB};
}

% =====================================================================
% 7. THE COUNTING
% =====================================================================
\art{
  \fill[void] (0,0) rectangle (16,9);
  \begin{scope}[blend mode=screen]
     \fill[ash,    opacity=0.35, path fading=rglow] (4.4,4.6) circle (5);
     \fill[silver, opacity=0.30, path fading=rglow] (11.6,2.4) circle (3.5);\end{scope}
  \dust{80}
  \begin{scope}[blend mode=screen]
     \fill[smoke, opacity=0.40, path fading=rglow] (4.4,4.6) ellipse (1.7 and 4.2);\end{scope}
  \shade[bottom color=slate, top color=silver] (3.5,0.9) rectangle (5.3,8.3);
  \draw[silver, line width=1pt, opacity=0.7] (3.5,0.9) rectangle (5.3,8.3);
  \node[text=ink, rotate=90, font=\footnotesize\bfseries] at (4.4,4.6) {11{,}695{,}232{,}963 comparisons};
  \spark{(11.6,2.4)}{0.14}{white}
  \node[text=silver, font=\footnotesize] at (11.6,3.4) {159{,}018 hash probes};
  \node[text=white] at (10.9,8.0) {\Large don't trust the number ---};
  \node[text=white] at (10.9,7.0) {\Large \textbf{count}};
  \node[text=ash, font=\footnotesize] at (10.9,5.6) {a list, walked \quad vs.\quad an index};
}

% =====================================================================
% 8. THE THREE ROTORS  (brightness = refinement)
% =====================================================================
\art{
  \fill[void] (0,0) rectangle (16,9);
  \begin{scope}[blend mode=screen]
     \fill[ash, opacity=0.28, path fading=rglow] (8,4.8) circle (9);\end{scope}
  \dust{70}
  \node[text=white] at (8,8.4) {\Large three rotors, one model --- the lesson travels};
  \draw[ash, line width=1.5pt, opacity=0.5, -{Stealth[length=4mm]}]
     (4.4,6.0) .. controls (6,7.0) and (7,7.0) .. (7.6,6.2);
  \draw[silver, line width=1.5pt, opacity=0.6, -{Stealth[length=4mm]}]
     (8.6,6.2) .. controls (10,7.0) and (11,7.0) .. (11.8,6.0);
  % dim -> bright orbs
  \fill[ash, opacity=0.30, path fading=rglow] (3.4,5.4) circle (1.5);
  \shade[inner color=smoke, outer color=slate] (3.4,5.4) circle (0.92);
  \draw[ash, line width=1pt] (3.4,5.4) circle (0.92);
  \fill[silver, opacity=0.40, path fading=rglow] (8.0,5.4) circle (1.6);
  \shade[inner color=white, outer color=ash] (8.0,5.4) circle (0.92);
  \draw[silver, line width=1.2pt] (8.0,5.4) circle (0.92);
  \rays{(12.6,5.4)}{white}{14}
  \fill[white, opacity=0.50, path fading=rglow] (12.6,5.4) circle (1.7);
  \shade[inner color=white, outer color=smoke] (12.6,5.4) circle (0.92);
  \draw[white, line width=1.4pt] (12.6,5.4) circle (0.92);
  % labels
  \node[text=smoke,  font=\small\bfseries] at (3.4,3.9) {Original C};
  \node[text=ash,    font=\footnotesize\itshape] at (3.4,3.5) {the master};
  \node[text=silver, font=\small\bfseries] at (8.0,3.9) {Hash-consed C};
  \node[text=smoke,  font=\footnotesize\itshape] at (8.0,3.5) {taught one trick};
  \node[text=white,  font=\small\bfseries] at (12.6,3.9) {Rust};
  \node[text=silver, font=\footnotesize\itshape] at (12.6,3.5) {the lesson, fully};
  \foreach \x in {3.4,8.0,12.6}{\draw[ash, opacity=0.4, line width=0.5pt] (\x-1.3,3.15)--(\x+1.3,3.15);}
  \node[text=smoke,  font=\footnotesize] at (3.4,2.7) {139\,s $\cdot$ 428\,MB};
  \node[text=ash,    font=\scriptsize]  at (3.4,2.3) {3.17\,M built $\cdot$ 11.7\,B compares};
  \node[text=silver, font=\footnotesize] at (8.0,2.7) {1.5\,s $\cdot$ byte-identical};
  \node[text=smoke,  font=\scriptsize]  at (8.0,2.3) {95$\times$ $\cdot$ 36/36 equal};
  \node[text=white,  font=\footnotesize] at (12.6,2.7) {0.1\,s $\cdot$ 20\,MB};
  \node[text=silver, font=\scriptsize]  at (12.6,2.3) {159\,k built $\cdot$ 1000$\times$};
  \node[text=silver, opacity=0.9] at (8,0.95)
     {\footnotesize the speed was never the language --- it was \textit{one small thing, done well}};
}

% =====================================================================
% 9. THE BACK-PORT
% =====================================================================
\art{
  \fill[void] (0,0) rectangle (16,9);
  \begin{scope}[blend mode=screen]
     \fill[ash,    opacity=0.30, path fading=rglow] (3.9,4.7) circle (3.4);
     \fill[silver, opacity=0.30, path fading=rglow] (12.1,4.7) circle (3.4);\end{scope}
  \dust{70}
  \rays{(3.9,4.7)}{smoke}{14}\rays{(12.1,4.7)}{white}{14}
  \shade[inner color=smoke, outer color=slate] (3.9,4.7) circle (1.05);
  \node[text=ink, font=\bfseries] at (3.9,4.7) {C};
  \shade[inner color=white, outer color=ash] (12.1,4.7) circle (1.05);
  \node[text=ink, font=\bfseries] at (12.1,4.7) {C\,+\,hash};
  \draw[white, line width=6pt, opacity=0.08] (5.0,4.7) -- (11.0,4.7);
  \draw[silver, line width=2.5pt, opacity=0.30,
        decorate, decoration={random steps, segment length=5pt, amplitude=2pt}] (5.0,4.7) -- (11.0,4.7);
  \draw[white, line width=1pt, opacity=0.9,
        decorate, decoration={random steps, segment length=5pt, amplitude=2pt}] (5.0,4.7) -- (11.0,4.7);
  \spark{(8,4.7)}{0.10}{white}
  \node[text=white] at (8,8.0) {\Large the one trick, carried back into the original itself};
  \glowtext{(8,2.2)}{silver}{\fontsize{38}{40}\selectfont\textbf{\textcolor{white}{95$\times$}} \normalsize\textcolor{silver}{, byte-identical}}
  \node[text=ash, font=\footnotesize] at (8,1.2) {``perfection is lots of little things done well''};
}

% =====================================================================
% 10. DEDUP OFF -- shatter vs serenity
% =====================================================================
\art{
  \fill[void] (0,0) rectangle (16,9);
  \begin{scope}[blend mode=screen]
     \fill[ash,    opacity=0.30, path fading=rglow] (4.3,4.6) circle (4);
     \fill[silver, opacity=0.30, path fading=rglow] (11.7,4.6) circle (4);\end{scope}
  \dust{60}
  % shattered crystal (dark fractured core)
  \rays{(4.3,4.6)}{smoke}{12}
  \shade[inner color=slate, outer color=ash]
     (4.3,6.5) -- (3.0,4.4) -- (3.8,2.4) -- (4.9,2.7) -- (5.6,4.7) -- cycle;
  \draw[white, opacity=0.9, line width=1pt] (4.3,6.5) -- (4.0,4.2);
  \draw[white, opacity=0.9, line width=1pt] (3.0,4.4) -- (4.4,3.7);
  \draw[white, opacity=0.9, line width=1pt] (5.6,4.7) -- (4.0,4.2);
  \draw[white, opacity=0.9, line width=1pt] (3.8,2.4) -- (4.6,4.4);
  \fill[slate, opacity=0.95] (6.1,5.6) -- (6.6,5.0) -- (6.3,4.5) -- cycle;
  \draw[ash, line width=0.6pt] (6.1,5.6) -- (6.6,5.0) -- (6.3,4.5) -- cycle;
  \node[text=white, font=\bfseries] at (4.3,1.4) {C: crashes};
  % serene gem (luminous core)
  \shade[inner color=white, outer color=ash]
     (11.7,6.4) -- (10.3,4.6) -- (11.7,2.7) -- (13.1,4.6) -- cycle;
  \draw[silver, opacity=0.8, line width=0.8pt]
     (11.7,6.4) -- (11.7,2.7) (10.3,4.6) -- (13.1,4.6)
     (11.7,6.4) -- (10.3,4.6) -- (11.7,2.7) -- (13.1,4.6) -- cycle;
  \begin{scope}[blend mode=screen]
     \fill[white, opacity=0.7, path fading=rglow] (11.7,5.0) circle (0.5);\end{scope}
  \node[text=white, font=\bfseries] at (11.7,1.4) {Rust: unchanged};
  \node[text=white] at (8,8.3) {\Large turn the sharing off, in both};
  \node[text=ash, font=\footnotesize] at (8,7.4) {it is load-bearing in C; merely cosmetic in Rust};
}

% =====================================================================
% 11. SYMBOLIC EXECUTION -- one beam over a sea of inputs
% =====================================================================
\art{
  \fill[void] (0,0) rectangle (16,9);
  \begin{scope}[blend mode=screen]
     \fill[ash, opacity=0.32, path fading=rglow] (8,7) circle (7);\end{scope}
  \foreach \x in {0,...,37}{\foreach \y in {0,...,12}{
     \pgfmathsetmacro{\jx}{\x*0.40+0.6}\pgfmathsetmacro{\jy}{\y*0.40+1.6}
     \fill[smoke, opacity={0.18+rnd*0.45}] (\jx,\jy) circle (0.05);}}
  \begin{scope}[blend mode=screen]
     \fill[silver, opacity=0.20, path fading=rglow] (8,8.4) -- (0.8,1.3) -- (15.2,1.3) -- cycle;
     \fill[white,  opacity=0.30, path fading=rglow] (8,8.4) -- (5.4,1.3) -- (10.6,1.3) -- cycle;
  \end{scope}
  \spark{(8,8.3)}{0.12}{white}
  \node[text=white] at (8,0.7) {\Large one question covers \textbf{every} input at once};
}

% =====================================================================
% 12. SYMBOLIC ARGV -- locks open
% =====================================================================
\art{
  \fill[void] (0,0) rectangle (16,9);
  \begin{scope}[blend mode=screen]
     \fill[ash, opacity=0.30, path fading=rglow] (8,4.6) circle (8);\end{scope}
  \dust{70}
  \foreach \i/\open in {0/0,1/0,2/1,3/1,4/1,5/1}{
     \pgfmathsetmacro{\lx}{2.0+\i*2.35}
     \ifnum\open=1
        \fill[white, opacity=0.40, path fading=rglow] (\lx,4.7) circle (1.5);
        \begin{scope}[blend mode=screen]
        \foreach \k in {1,...,7}{
           \pgfmathsetmacro{\dy}{rnd*2.4}\pgfmathsetmacro{\dx}{(rnd-0.5)*0.7}
           \fill[white, opacity={0.5-0.15*rnd}, path fading=rglow] (\lx+\dx,5.4+\dy) circle (0.07);}
        \end{scope}
        \def\lc{white}
     \else \def\lc{ash}\fi
     \shade[top color=\lc!75!void, bottom color=\lc!30!void] (\lx-0.5,3.5) rectangle (\lx+0.5,4.7);
     \draw[\lc, line width=1.5pt] (\lx-0.5,3.5) rectangle (\lx+0.5,4.7);
     \ifnum\open=1 \draw[\lc, line width=1.6pt] (\lx-0.30,4.7) arc (200:20:0.36);
     \else \draw[\lc, line width=1.6pt] (\lx-0.30,4.7) arc (180:0:0.30);\fi
     \fill[\lc] (\lx,4.1) circle (0.08);}
  \node[text=white] at (8,8.2) {\Large the command line: \textcolor{ash}{frozen} $\rightarrow$ \textbf{open}};
  \node[text=silver, font=\normalsize] at (8,2.0) {the bytes after the program's name become \textit{questions}};
}

% =====================================================================
% 13. THE X / Y SECRET
% =====================================================================
\art{
  \fill[void] (0,0) rectangle (16,9);
  \begin{scope}[blend mode=screen]
     \fill[ash, opacity=0.34, path fading=rglow] (8,5.0) circle (8);\end{scope}
  \dust{70}
  \foreach \dx/\ch in {5.4/X, 10.6/Y}{
     \fill[silver, opacity=0.35, path fading=rglow] (\dx,5.1) circle (2.0);
     \shade[inner color=slate, outer color=void] (\dx,5.1) circle (1.35);
     \draw[white, line width=2pt] (\dx,5.1) circle (1.35);
     \draw[silver, line width=0.8pt, opacity=0.6] (\dx,5.1) circle (1.05);
     \foreach \t in {0,15,...,345}{\draw[silver,opacity=0.6,line width=1pt]
        ({\dx+1.16*cos(\t)},{5.1+1.16*sin(\t)}) -- ({\dx+1.35*cos(\t)},{5.1+1.35*sin(\t)});}
     \node[text=white] at (\dx,5.1) {\fontsize{40}{42}\selectfont\textbf{\ch}};}
  \rays{(8,5.1)}{white}{30}
  \begin{scope}[blend mode=screen]
  \foreach \t in {0,18,...,342}{
     \pgfmathsetmacro{\rr}{0.7+0.6*rnd}
     \draw[white,opacity=0.8,line width=1.2pt]
        ({8+0.3*cos(\t)},{5.1+0.3*sin(\t)}) -- ({8+\rr*cos(\t)},{5.1+\rr*sin(\t)});}
  \end{scope}
  \spark{(8,5.1)}{0.10}{white}
  \node[text=white] at (8,8.5) {\Large it crashes only on \textbf{X} and \textbf{Y} together};
  \node[text=silver] at (8,1.2) {\large no keyboard could trigger it --- the solver found it by reason};
}

% =====================================================================
% 14. DIFFERENTIAL VALIDATION -- 36 / 36
% =====================================================================
\art{
  \fill[void] (0,0) rectangle (16,9);
  \begin{scope}[blend mode=screen]
     \fill[silver, opacity=0.26, path fading=rglow] (8,5.4) circle (7);\end{scope}
  \dust{60}\rays{(8,5.4)}{silver}{16}
  \foreach \r in {0,...,3}{\foreach \cc in {0,...,8}{
     \pgfmathsetmacro{\cx}{2.6+\cc*1.35}
     \pgfmathsetmacro{\cy}{7.0-\r*1.05}
     \fill[silver, opacity=0.45, path fading=rglow] (\cx,\cy) circle (0.42);
     \draw[white, line width=1.8pt, line cap=round]
        (\cx-0.18,\cy) -- (\cx-0.03,\cy-0.16) -- (\cx+0.24,\cy+0.22);}}
  \glowtext{(8,1.7)}{silver}{\fontsize{46}{48}\selectfont\textbf{\textcolor{white}{36 / 36}}}
  \node[text=ash, font=\footnotesize] at (8,0.75)
    {same bug, same step --- both tools, every benchmark, two configurations};
}

% =====================================================================
% 15. THE VISUALIZER
% =====================================================================
\art{
  \fill[void] (0,0) rectangle (16,9);
  \begin{scope}[blend mode=screen]
     \fill[ash, opacity=0.30, path fading=rglow] (8,4.4) circle (8);\end{scope}
  \dust{70}
  \coordinate (top) at (8,6.7);
  \coordinate (na) at (5.6,5.4); \coordinate (nb) at (10.4,5.4);
  \coordinate (nc) at (3.4,3.6); \coordinate (nd) at (6.8,3.6);
  \coordinate (ne) at (9.2,3.6); \coordinate (nf) at (12.2,3.6);
  \coordinate (ng) at (4.8,1.9); \coordinate (nh) at (7.6,1.9);
  \foreach \a/\b in {top/na, top/nb, na/nc, na/nd, nb/ne, nb/nf, nd/ng, nd/nh}{\neon{(\a)}{(\b)}{silver}}
  % bad state: an inverted node -- white-filled star, black text, strong halo
  \rays{(8,7.4)}{white}{20}
  \fill[white, opacity=0.55, path fading=rglow] (8,7.4) circle (1.6);
  \node[star, star points=8, star point ratio=0.55, draw=white, fill=void,
        line width=1.6pt, minimum size=1.7cm] at (8,7.4) {};
  \node[text=white, font=\tiny\bfseries] at (8,7.4) {VIOLATED};
  \foreach \n in {na,nb,nc,nd,ne,nf,ng,nh}{\spark{(\n)}{0.085}{silver}}
  \node[text=white] at (8,0.8) {\Large the solver's wall of binary, made a picture you can point at};
}

% =====================================================================
% 16. CLOSE
% =====================================================================
\art{
  \fill[void] (0,0) rectangle (16,9);
  \mist \dust{230}\bokehm{45}
  \rings{(8,5.2)}{7}{1.05}{ash}
  \coordinate (r1) at (4.6,6.3);\coordinate (r2) at (8.2,7.0);\coordinate (r3) at (11.4,5.6);
  \neon{(r1)}{(r2)}{silver}\neon{(r2)}{(r3)}{silver}
  \spark{(r1)}{0.09}{white}\spark{(r2)}{0.09}{white}\spark{(r3)}{0.09}{white}
  \begin{scope}[blend mode=screen]
     \fill[silver, opacity=0.40, path fading=rglow] (8,3.55) ellipse (6.2 and 1.4);\end{scope}
  \node[text=white] at (8,3.6) {\fontsize{28}{32}\selectfont\textbf{proving it still tells the truth}};
  \node[text=silver, opacity=0.9] at (8,2.6) {\large was the whole craft. Thank you.};
  \node[text=ash] at (8,1.4)
    {\footnotesize github.com/jaspek/rotor-rust \;\textbar\; rotor-c-hashcons
     \;\textbar\; jaspek.github.io/rotor-rust};
}

\end{document}
